\documentclass[../LuanVan.tex]{subfiles}
\begin{document}

\section{Thực nghiệm so sánh các phương pháp phát hiện khuôn mặt}
\label{section:5.1}

\subsection{Tập dữ liệu và giao thức đánh giá}
\label{subsection:5.1.1}

Thực nghiệm sử dụng toàn bộ tập xác thực của WIDER FACE gồm 3.226 ảnh. WIDER FACE được xây dựng cho bài toán phát hiện khuôn mặt; toàn bộ tập dữ liệu chứa 32.203 ảnh với 393.703 khung bao khuôn mặt và có biến thiên lớn về tỷ lệ, tư thế, mức độ che khuất, biểu cảm và điều kiện chiếu sáng \cite{yang2016wider}. Tập xác thực được sử dụng vì có nhãn chuẩn được công bố: vị trí khuôn mặt trên mỗi ảnh đã được đánh dấu bằng khung bao và dữ liệu được chia thành ba nhóm độ khó gồm dễ, trung bình và khó. Các đặc điểm này cho phép đánh giá ba mô hình phát hiện trên cùng một dữ liệu và cùng tiêu chuẩn. Tập dữ liệu FER gồm các ảnh đã cắt sát khuôn mặt không thể thay thế tập dữ liệu chuyên dùng để đánh giá phát hiện khuôn mặt.

Cấu hình của ba mô hình phát hiện được cố định trước khi đo. Haar Cascade sử dụng tệp \path{haarcascade_frontalface_default.xml} của OpenCV \cite{opencvHaar}. MTCNN sử dụng mô hình tiền huấn luyện của \texttt{facenet-pytorch} phiên bản 2.6.0 \cite{facenetPytorch}, với \texttt{min\_face\_size=20}, các ngưỡng $[0{,}6;0{,}7;0{,}7]$ và hệ số tỷ lệ $0{,}709$. RetinaFace sử dụng mạng nền MobileNet0.25 cùng trọng số tiền huấn luyện trên tập huấn luyện WIDER FACE từ kho mã \path{Pytorch_Retinaface} \cite{pytorchRetinaface}. Việc cố định một cấu hình cho mỗi phương pháp giúp kết quả giữa các lượt chạy có thể tái lập và tránh trộn lẫn ảnh hưởng của nhiều mạng nền hoặc thư viện triển khai.

Ba mô hình nhận cùng ảnh gốc và tự thay đổi kích thước ảnh theo quy trình riêng. Trước khi đo chính thức, mỗi mô hình xử lý thử 20 ảnh; kết quả của các ảnh này không được đưa vào số liệu báo cáo. Bước chạy thử giúp hoàn tất việc khởi tạo bộ nhớ và các thành phần tính toán, qua đó giảm ảnh hưởng bất thường của những lượt suy luận đầu tiên đến kết quả đo thời gian.

Sau đó, mỗi mô hình được chạy ba lượt trên cùng danh sách ảnh và thứ tự mô hình được xoay vòng giữa các lượt. Thời gian đo bao gồm các bước tiền xử lý, suy luận và hậu xử lý của mô hình phát hiện khuôn mặt, nhưng không bao gồm thời gian đọc ảnh, ghi tệp hoặc tính các chỉ số đánh giá.

Thực nghiệm sử dụng CPU với một luồng PyTorch và một luồng OpenCV trên Apple M4 Pro, RAM 24 GiB và macOS 26.5. Môi trường phần mềm gồm Python 3.11.7, OpenCV 4.10.0 và PyTorch 2.2.2.

Chất lượng chính được đánh giá bằng độ chính xác trung bình (AP) trên ba mức khó theo giao thức WIDER FACE, với ngưỡng tỷ lệ phần giao trên phần hợp (IoU) bằng 0,5. Với khung bao dự đoán $B_p$ và khung bao chuẩn $B_g$, IoU được xác định bởi

\begin{equation}
    \operatorname{IoU}(B_p,B_g)=\frac{|B_p\cap B_g|}{|B_p\cup B_g|}.
    \label{eq:iou}
\end{equation}

Mỗi khung bao dự đoán chỉ được ghép với một khung bao chuẩn. Gọi $P$ là độ chính xác và $R$ là độ bao phủ tại ngưỡng đang xét; các chỉ số $P$, $R$ và F1 được tính theo

\begin{equation}
    P=\frac{TP}{TP+FP},\quad
    R=\frac{TP}{TP+FN},\quad
    F1=\frac{2PR}{P+R}.
    \label{eq:detection_metrics}
\end{equation}

Ngưỡng IoU bảo đảm một dự đoán chỉ được tính là đúng khi khung bao không những chỉ ra sự tồn tại của khuôn mặt mà còn định vị đủ gần với khung chuẩn. Một dự đoán thỏa ngưỡng và được ghép với một khung chuẩn chưa được sử dụng được tính là dương tính đúng ($TP$). Khung dự đoán không ghép được tính là dương tính sai ($FP$), còn khuôn mặt trong nhãn chuẩn không có dự đoán tương ứng được tính là âm tính sai ($FN$).

Độ chính xác $P$ cho biết trong các vùng được mô hình dự đoán là khuôn mặt, tỷ lệ dự đoán đúng là bao nhiêu; chỉ số này phản ánh mức độ phát hiện nhầm nền hoặc vật thể khác thành khuôn mặt. Độ bao phủ $R$ cho biết trong toàn bộ khuôn mặt có trong dữ liệu, mô hình phát hiện được bao nhiêu; chỉ số này phản ánh mức độ bỏ sót. Độ bao phủ có ý nghĩa trực tiếp với hệ thống đề xuất vì một khuôn mặt bị bỏ sót sẽ không tạo được vùng đầu vào cho các bước phân loại biểu cảm và tạo đặc trưng nhận dạng. Tuy nhiên, chỉ sử dụng riêng một trong hai chỉ số có thể dẫn đến đánh giá lệch: giảm ngưỡng tin cậy thường làm tăng độ bao phủ nhưng cũng có thể tăng số phát hiện nhầm, trong khi tăng ngưỡng có thể cải thiện độ chính xác nhưng làm bỏ sót nhiều khuôn mặt hơn. F1 là trung bình điều hòa của $P$ và $R$, được sử dụng để mô tả mức cân bằng giữa hai loại sai sót tại một ngưỡng vận hành xác định.

AP là diện tích dưới đường độ chính xác--độ bao phủ sau nội suy theo giao thức đánh giá đối tượng \cite{everingham2010pascal,yang2016wider}. Khác với $P$, $R$ và F1 được tính tại một ngưỡng cố định, AP tổng hợp kết quả khi thay đổi ngưỡng tin cậy, qua đó phản ánh khả năng xếp hạng các dự đoán đúng cao hơn các dự đoán sai trên toàn bộ dải ngưỡng. AP được chọn làm chỉ số chất lượng chính vì mức tin cậy của ba mô hình không được hiệu chỉnh trên cùng một thang xác suất; việc chỉ so sánh tại một giá trị ngưỡng có thể tạo lợi thế cho một cấu hình cụ thể. Các chỉ số $P$, $R$ và F1 vẫn được báo cáo để mô tả hành vi của từng mô hình tại ngưỡng triển khai đã cố định: Haar Cascade không lọc thêm, MTCNN dùng đầu ra sau ngưỡng O-Net 0,7 và RetinaFace dùng ngưỡng tin cậy 0,6.

AP được trình bày riêng cho ba nhóm dễ, trung bình và khó nhằm tránh trường hợp kết quả tốt trên các khuôn mặt lớn, rõ che khuất việc mô hình bỏ sót khuôn mặt nhỏ, bị che hoặc xuất hiện trong cảnh đông người. Cách phân tách này cho phép quan sát mức suy giảm của từng phương pháp khi độ khó tăng và xác định phương pháp duy trì chất lượng ổn định hơn trong điều kiện phức tạp. Như vậy, AP trên ba mức khó cung cấp đánh giá chất lượng tổng thể, còn $P$, $R$ và F1 giải thích cụ thể hai loại sai sót tại ngưỡng vận hành. Hiệu năng tính toán được đánh giá riêng bằng độ trễ trung vị, độ trễ tại phân vị thứ 95 và tốc độ xử lý tính theo số khung hình mỗi giây (FPS), vì một mô hình phát hiện tốt nhưng quá chậm vẫn có thể không phù hợp với luồng xử lý của hệ thống.

\subsection{Kết quả thực nghiệm so sánh}
\label{subsection:5.1.2}

Bảng \ref{tab:detector_quality} và Bảng \ref{tab:detector_performance} được tạo từ cùng lần chạy \path{wider-val-cpu-20260921}. Mỗi mô hình có 9.678 bản ghi thời gian, tương ứng với 3.226 ảnh qua ba lượt chạy. Kết quả AP đã được đối chiếu bằng bộ đánh giá WIDER FACE độc lập và trùng khớp đến sai số làm tròn.

\begin{table}[H]
    \centering
    \caption{Kết quả chất lượng phát hiện trên tập xác thực WIDER FACE}
    \label{tab:detector_quality}
    \resizebox{\textwidth}{!}{%
    \begin{tabular}{|l|r|r|r|r|r|r|}
        \hline
        \textbf{Mô hình} & \textbf{AP dễ} & \textbf{AP trung bình} & \textbf{AP khó} & \textbf{Độ chính xác} & \textbf{Độ bao phủ} & \textbf{F1} \\
        \hline
        Haar Cascade -- OpenCV, cấu hình mặt chính diện mặc định & 0,4252 & 0,3668 & 0,1754 & 0,6906 & 0,1514 & 0,2484 \\
        \hline
        MTCNN -- facenet-pytorch & 0,8132 & 0,7706 & 0,4990 & 0,9081 & 0,4063 & 0,5614 \\
        \hline
        RetinaFace -- MobileNet0.25 & \textbf{0,9082} & \textbf{0,8830} & \textbf{0,7388} & \textbf{0,9106} & \textbf{0,4784} & \textbf{0,6273} \\
        \hline
    \end{tabular}}
\end{table}

\begin{table}[H]
    \centering
    \caption{Kết quả hiệu năng phát hiện trên CPU}
    \label{tab:detector_performance}
    \resizebox{\textwidth}{!}{%
    \begin{tabular}{|l|r|r|r|r|}
        \hline
        \textbf{Mô hình} & \textbf{Trung vị (ms)} & \textbf{Phân vị 95 (ms)} & \textbf{Số ảnh/giây} & \textbf{Tỷ lệ ảnh không phát hiện} \\
        \hline
        Haar Cascade & 131,52 & 257,71 & 7,03 & 0,2228 \\
        \hline
        MTCNN & 197,95 & 418,26 & 4,55 & 0,0615 \\
        \hline
        RetinaFace--MobileNet0.25 & \textbf{64,53} & \textbf{131,99} & \textbf{13,01} & \textbf{0,0320} \\
        \hline
    \end{tabular}}
\end{table}

\subsection{Nhận xét từ kết quả thực nghiệm}
\label{subsection:5.1.3}

Kết quả thực nghiệm cho thấy ba phương pháp khác nhau cả về chất lượng phát hiện và chi phí suy luận. Haar Cascade có kết quả thấp nhất trên cả ba nhóm dễ, trung bình và khó. Mức giảm mạnh ở nhóm khó cho thấy cấu hình phát hiện mặt chính diện mặc định gặp khó khăn khi dữ liệu có nhiều khuôn mặt nhỏ, bị che khuất hoặc không ở tư thế chính diện. MTCNN cải thiện rõ rệt AP và giảm tỷ lệ ảnh không phát hiện so với Haar Cascade. Đổi lại, MTCNN có độ trễ trung vị và độ trễ tại phân vị thứ 95 cao nhất, phản ánh chi phí của quá trình xử lý nối tiếp qua ba tầng trên tháp ảnh đa tỷ lệ trong cấu hình thực nghiệm.

RetinaFace sử dụng mạng nền MobileNet0.25 duy trì AP cao nhất ở cả ba mức khó, đồng thời đạt độ trễ thấp nhất và số ảnh xử lý trong một giây cao nhất. Khoảng cách 0,2398 AP so với MTCNN trên nhóm khó cho thấy lợi thế của RetinaFace rõ hơn khi điều kiện phát hiện trở nên phức tạp. Tỷ lệ ảnh không phát hiện 0,0320 cũng cho thấy số ảnh không tạo được đầu vào cho bước phân loại biểu cảm thấp hơn so với hai phương pháp còn lại. Trong phạm vi tập dữ liệu WIDER FACE, trên thiết bị CPU và với bộ tham số đã cố định, RetinaFace sử dụng mạng nền MobileNet0.25 thể hiện sự cân bằng tốt nhất giữa chất lượng phát hiện và tốc độ xử lý. Kết quả này chỉ phản ánh cấu hình thực nghiệm đã mô tả, không phải kết luận khái quát cho mọi dữ liệu hoặc nền tảng phần cứng.

\section{Thực nghiệm vận hành hệ thống}
\label{section:5.2}

\subsection{Mục tiêu và tiêu chí kiểm nghiệm}
\label{subsection:5.2.1}

Thực nghiệm vận hành nhằm trả lời câu hỏi liệu các thành phần đã xây dựng có phối hợp được thành một hệ thống hoàn chỉnh hay không. Phạm vi kiểm nghiệm bắt đầu từ ảnh đầu vào, đi qua phát hiện khuôn mặt, phân loại biểu cảm, tạo và đối sánh véc-tơ ArcFace, hình thành quan sát và lần mua sắm, sau đó kết thúc ở đơn hàng và báo cáo CRM. Đây là kiểm nghiệm tính hoạt động trong điều kiện dữ liệu được kiểm soát, không phải phép đo độ chính xác ngoài thực tế của nhận dạng danh tính hoặc phân loại biểu cảm.

Các tiêu chí được xác định trước khi chạy và trình bày tại Bảng~\ref{tab:operational_criteria_ch5}. Một tiêu chí chỉ được tính là đạt khi kết quả trong phản hồi của hệ thống đồng thời khớp với dữ liệu lưu trong PostgreSQL.

\begin{table}[H]
    \centering
    \caption{Tiêu chí kiểm nghiệm tính hoạt động của hệ thống}
    \label{tab:operational_criteria_ch5}
    \resizebox{\textwidth}{!}{%
    \begin{tabular}{|l|p{12.5cm}|}
        \hline
        \textbf{Mã} & \textbf{Điều kiện đạt} \\
        \hline
        C1 & Tạo đủ 100 hồ sơ khách hàng và 100 mẫu khuôn mặt từ ảnh đăng ký. \\
        \hline
        C2 & Mỗi ảnh quan sát của khách đã đăng ký tạo đúng một quan sát và được liên kết với đúng mã khách hàng mong đợi. \\
        \hline
        C3 & Ảnh của người không thuộc tập đăng ký được giữ ở trạng thái chưa xác định, không bị gắn với hồ sơ gần nhất. \\
        \hline
        C4 & Khung hình có hai khuôn mặt tạo hai quan sát; ảnh không có khuôn mặt và dữ liệu không thể giải mã không tạo quan sát giả. \\
        \hline
        C5 & Các quan sát hình thành được lần mua sắm, liên kết được đơn hàng và tạo dữ liệu cho các báo cáo phân bố, biến thiên, thay đổi giữa khu vực và chất lượng dữ liệu. \\
        \hline
        C6 & Các dịch vụ đạt trạng thái sẵn sàng; kiểm thử máy chủ, dịch vụ xử lý ảnh và quá trình biên dịch giao diện đều hoàn thành. \\
        \hline
    \end{tabular}}
\end{table}

\subsection{Xây dựng dữ liệu thử nghiệm và lý do lựa chọn}
\label{subsection:5.2.2}

Trong phạm vi đề tài, điều kiện triển khai chưa cho phép lắp đặt đồng bộ camera tại các khu vực của cửa hàng và kết nối với hệ thống bán hàng thực tế. Vì vậy, dữ liệu dùng cho thực nghiệm được xây dựng theo một kịch bản có kiểm soát. Mục đích của cách xây dựng này là cung cấp đủ đầu vào để kiểm tra luồng xử lý ảnh, quản lý khách hàng, hình thành lần mua sắm, liên kết đơn hàng và lập báo cáo; dữ liệu không được xem là bản ghi hoạt động của một cửa hàng đang vận hành.

Nguồn ảnh duy nhất được sử dụng là tập xác thực FairFace. FairFace được công bố cho nghiên cứu thuộc tính khuôn mặt và có giấy phép CC BY 4.0 \cite{karkkainen2021fairface,fairfaceRepository}. Từ tệp nguồn, 110 bản ghi được chọn và chia thành ba phần không giao nhau: 100 bản ghi đại diện cho 100 khách hàng đã đăng ký, tám bản ghi dùng để xác định ngưỡng đối sánh và hai bản ghi dùng để kiểm tra trường hợp người chưa đăng ký. Chỉ số nguồn và mã kiểm tra SHA-256 của từng ảnh được lưu trong tệp kê khai để có thể truy vết. Các nhãn giới tính, nhóm tuổi và chủng tộc của FairFace không được đưa vào hệ thống.

Mỗi bản ghi thuộc nhóm khách hàng được dùng để xây dựng ba loại ảnh. Ảnh hồ sơ có kích thước $256\times256$ pixel và chỉ phục vụ hiển thị trên CRM. Khung đăng ký có kích thước $640\times480$ pixel; RetinaFace xác định vùng khuôn mặt và ArcFace chuyển vùng này thành véc-tơ 512 chiều để lưu làm mẫu đối sánh. Khung quan sát có cùng độ phân giải nhưng khuôn mặt được thay đổi vị trí, kích thước và độ sáng nhằm tạo sự khác biệt giữa thời điểm đăng ký và thời điểm xuất hiện tại cửa hàng. Vì vậy, kết quả nhận dạng được tính từ ảnh quan sát qua mô hình, không được lấy trực tiếp từ ảnh hồ sơ hoặc gán trước theo mã khách hàng.

Dữ liệu từ camera được thay thế bằng các khung ảnh có thông tin khu vực và thời gian ghi nhận. Năm khách hàng đầu có bốn khung quan sát tương ứng với bốn khu vực; 95 khách hàng còn lại có một khung quan sát, tạo thành 115 sự kiện của khách đã đăng ký. Trường hợp nhiều người trong cùng khung hình được xây dựng bằng cách đặt hai ảnh FairFace khác nhau vào một sự kiện thu nhận. Sự kiện này có một mã khung hình nhưng phải tạo hai quan sát, mỗi quan sát có chỉ số khuôn mặt, khung bao, kết quả biểu cảm và trạng thái nhận dạng riêng. Cách tổ chức này kiểm tra việc hệ thống không loại bỏ cả khung hình khi phát hiện nhiều hơn một người.

Trước lần chạy chính thức, ảnh của nhóm khách hàng được sàng lọc để bảo đảm RetinaFace tạo được vùng khuôn mặt trong cả khung đăng ký và khung quan sát. Đồng thời, véc-tơ của ảnh quan sát phải gần véc-tơ đăng ký cùng nguồn hơn các mẫu khác dưới ngưỡng đã xác định. Một bản ghi không đáp ứng điều kiện này được thay bằng bản ghi kế tiếp chưa sử dụng. Tám bản ghi dành cho xác định ngưỡng, mỗi bản ghi có ba biến thể, cung cấp 24 cặp cùng nguồn và 252 cặp khác nguồn. Khoảng cách cosine được tính trên véc-tơ ArcFace và ngưỡng thu được là 0,269958. Bước sàng lọc giúp tạo đầu vào có kết quả mong đợi để kiểm tra luồng phần mềm; vì vậy, kết quả không được diễn giải thành độ chính xác nhận dạng trên toàn bộ FairFace.

Hai bản ghi không thuộc nhóm khách hàng chỉ xuất hiện ở pha quan sát để kiểm tra trường hợp không tìm thấy mẫu phù hợp. Ba trường hợp đầu vào bổ sung gồm một khung không có khuôn mặt, một tệp không thể giải mã thành ảnh và một khung chứa hai người. Tổng cộng, lớp xử lý ảnh có 120 sự kiện: 115 sự kiện của khách đã đăng ký, hai sự kiện của người chưa đăng ký và ba sự kiện kiểm tra điều kiện biên. Tất cả các sự kiện này được gửi qua cùng giao diện tiếp nhận dành cho nguồn camera.

Do chưa có luồng camera hoạt động liên tục và chưa kết nối được dữ liệu bán hàng, một bộ dữ liệu kiểm tra nghiệp vụ được xây dựng theo quy tắc cố định để tạo đủ quan hệ cho các chức năng CRM và báo cáo. Bộ dữ liệu sử dụng chính 100 khách hàng nêu trên, bốn khu vực, sáu nhóm sản phẩm và 24 sản phẩm; không bổ sung nguồn ảnh khuôn mặt thứ hai. Với hạt giống \texttt{20260923}, kịch bản tạo 125 lần mua sắm, trong đó mỗi lần đi qua từ hai đến bốn khu vực, cùng 90 đơn hàng mới, 12 quan sát chưa xác định và 12 sự kiện đầu vào không hợp lệ. Các nhãn biểu cảm trong phần dữ liệu này được gán theo quy tắc cố định để kiểm tra phép lọc, phép nhóm và cách trình bày báo cáo; chúng không đi qua DeepFace và được tách khỏi kết quả đánh giá luồng xử lý ảnh.

Hình~\ref{fig:data_construction_ch5} minh họa quan hệ giữa ảnh nguồn, khung đăng ký và khung quan sát. Toàn bộ quy tắc lựa chọn ảnh, biến đổi khung hình, thời điểm và khu vực đều được cố định để lần chạy có thể được thực hiện lại và đối chiếu theo từng mã sự kiện.

\begin{figure}[H]
    \centering
    \begin{minipage}[t]{0.30\textwidth}
        \centering
        \includegraphics[width=\textwidth]{Hinhve/Chuong5/5_01_anh_ho_so.jpg}\\
        \small Ảnh hồ sơ
    \end{minipage}\hfill
    \begin{minipage}[t]{0.32\textwidth}
        \centering
        \includegraphics[width=\textwidth]{Hinhve/Chuong5/5_02_khung_dang_ky.jpg}\\
        \small Khung đăng ký
    \end{minipage}\hfill
    \begin{minipage}[t]{0.32\textwidth}
        \centering
        \includegraphics[width=\textwidth]{Hinhve/Chuong5/5_03_khung_quan_sat.jpg}\\
        \small Khung quan sát
    \end{minipage}
    \caption{Ba dạng dữ liệu được tạo từ cùng một bản ghi FairFace}
    \label{fig:data_construction_ch5}
\end{figure}

\subsection{Quy trình chạy thực nghiệm}
\label{subsection:5.2.3}

Quy trình được thực hiện theo hai lớp để chỉ rõ thành phần nào đang được kiểm tra. Ở lớp xử lý ảnh, các dịch vụ được khởi động với RetinaFace--MobileNet0.25, mô hình Emotion qua DeepFace và ArcFace. Trước khi thực nghiệm, cơ sở dữ liệu được điều chỉnh để một ảnh camera có thể chứa nhiều khuôn mặt và kết quả của từng khuôn mặt được lưu thành một quan sát riêng. Hệ thống cũng lưu các trường hợp ảnh lỗi và lịch sử chỉnh sửa kết quả quan sát. Sau khi nạp 100 hồ sơ, hệ thống tạo 100 mẫu khuôn mặt, nạp ngưỡng đối sánh và lần lượt gửi 120 sự kiện ảnh qua cùng giao diện tiếp nhận dành cho nguồn camera. Với mỗi sự kiện, bộ kiểm tra đối chiếu trạng thái ảnh, số khuôn mặt, trạng thái nhận dạng và mã khách hàng thực tế với kết quả mong đợi.

Ở lớp nghiệp vụ, bộ dữ liệu đã mô tả tại Mục~\ref{subsection:5.2.2} được nạp qua giao diện dành riêng cho kiểm tra hệ thống. Các bản ghi được đánh dấu bằng nguồn dữ liệu và mã lần chạy để tách khỏi 120 sự kiện đã đi qua dịch vụ xử lý ảnh. Bộ kiểm tra sau khi nạp đối chiếu số lần mua sắm, chuỗi khu vực, liên kết đơn hàng và kết quả của các truy vấn báo cáo. Các bản ghi thuộc lớp này chỉ được dùng để chứng minh chức năng CRM và báo cáo, không được cộng vào kết quả C2--C4 hoặc dùng làm bằng chứng về chất lượng của RetinaFace, DeepFace và ArcFace.

\subsection{Kết quả của luồng xử lý ảnh}
\label{subsection:5.2.4}

Kết quả phát lại đạt 120/120 sự kiện theo các điều kiện đã xác định. Chi tiết được trình bày tại Bảng~\ref{tab:image_pipeline_results_ch5}.

\begin{table}[H]
    \centering
    \caption{Kết quả kiểm nghiệm luồng xử lý ảnh}
    \label{tab:image_pipeline_results_ch5}
    \resizebox{\textwidth}{!}{%
    \begin{tabular}{|l|r|r|l|}
        \hline
        \textbf{Nhóm đầu vào} & \textbf{Sự kiện} & \textbf{Quan sát tạo ra} & \textbf{Kết quả} \\
        \hline
        Khách đã đăng ký & 115 & 115 & Phát hiện một khuôn mặt, phân loại hợp lệ và liên kết đúng khách hàng \\
        \hline
        Người chưa đăng ký & 2 & 2 & Cả hai giữ trạng thái \texttt{NO\_MATCH} \\
        \hline
        Khung hình có hai người & 1 & 2 & Hai khuôn mặt được xử lý độc lập; cả hai giữ trạng thái \texttt{NO\_MATCH} \\
        \hline
        Khung hình không có khuôn mặt & 1 & 0 & Sự kiện mang trạng thái \texttt{NO\_FACE} \\
        \hline
        Dữ liệu ảnh không hợp lệ & 1 & 0 & Sự kiện mang trạng thái \texttt{INVALID\_IMAGE} \\
        \hline
        \textbf{Tổng} & \textbf{120} & \textbf{119} & \textbf{120/120 sự kiện đạt} \\
        \hline
    \end{tabular}}
\end{table}

Toàn bộ 119 khuôn mặt phát hiện được đều có kết quả phân loại biểu cảm hợp lệ. Trong đó 115 quan sát được đối sánh với đúng mã khách hàng; bốn quan sát còn lại gồm hai ảnh người chưa đăng ký và hai khuôn mặt trong cùng một khung hình nên không tạo chuỗi cá nhân. Kết quả này xác nhận hệ thống không yêu cầu khung hình chỉ có một người: một sự kiện thu nhận có thể sinh nhiều quan sát, còn ảnh không có khuôn mặt dừng ở cấp sự kiện.

Hình~\ref{fig:face_data_result_ch5} thể hiện danh sách khách hàng sau bước ghi nhận đồng ý và đăng ký mẫu khuôn mặt. Hình~\ref{fig:unidentified_result_ch5} cho thấy bản ghi chưa xác định vẫn có kết quả biểu cảm hợp lệ thay vì trạng thái \texttt{NOT\_RUN}; các trạng thái không chạy chỉ xuất hiện đúng ở ca không có khuôn mặt hoặc ảnh không thể giải mã.

\begin{figure}[H]
    \centering
    \includegraphics[width=0.98\textwidth]{Hinhve/Chuong5/5_04_du_lieu_khuon_mat.png}
    \caption{Danh sách khách hàng sau bước đăng ký dữ liệu khuôn mặt}
    \label{fig:face_data_result_ch5}
\end{figure}

\begin{figure}[H]
    \centering
    \includegraphics[width=0.98\textwidth]{Hinhve/Chuong5/5_05_quan_sat_chua_xac_dinh.png}
    \caption{Quan sát chưa xác định vẫn có kết quả phân loại biểu cảm hợp lệ}
    \label{fig:unidentified_result_ch5}
\end{figure}

\subsection{Kết quả tích hợp CRM và báo cáo}
\label{subsection:5.2.5}

Sau khi hoàn tất cả hai lớp, PostgreSQL có 100 khách hàng, 100 mẫu khuôn mặt, 529 sự kiện thu nhận, 516 quan sát, 225 lần mua sắm, 93 đơn hàng và 232 dòng sản phẩm. Trong 529 sự kiện, 120 sự kiện thuộc lớp xử lý ảnh; 409 sự kiện thuộc tải nghiệp vụ, gồm 397 sự kiện hợp lệ và 12 sự kiện lỗi. Trong 516 quan sát, 119 quan sát do mô hình xử lý ảnh tạo ra và 397 quan sát thuộc tải nghiệp vụ. Sự phân tách này được giữ bằng trường nguồn dữ liệu và mã lần chạy.

\begin{table}[H]
    \centering
    \caption{Số lượng dữ liệu sau thực nghiệm tích hợp}
    \label{tab:integrated_results_ch5}
    \begin{tabular}{|l|r|}
        \hline
        \textbf{Loại dữ liệu} & \textbf{Số lượng} \\
        \hline
        Khách hàng / mẫu khuôn mặt & 100 / 100 \\
        \hline
        Sự kiện thu nhận / quan sát & 529 / 516 \\
        \hline
        Lần mua sắm & 225 \\
        \hline
        Nhóm sản phẩm / sản phẩm & 6 / 24 \\
        \hline
        Đơn hàng / dòng sản phẩm & 93 / 232 \\
        \hline
    \end{tabular}
\end{table}

Tải nghiệp vụ tạo đúng 125 lần mua sắm và 90 đơn hàng mới; ba đơn còn lại là dữ liệu lịch sử được nạp trước để kiểm tra việc giữ nguyên tên và giá sản phẩm. Lớp xử lý ảnh tạo thêm 100 lần mua sắm từ 100 khách hàng. Bộ kiểm tra hậu điều kiện đánh giá 105 điều kiện, gồm mỗi khách hàng có ít nhất một lần mua sắm và năm khách hàng đầu có đủ chuỗi bốn khu vực; kết quả đạt 105/105. Các truy vấn báo cáo phân bố và thay đổi giữa khu vực đều trả về dữ liệu không rỗng.

Hình~\ref{fig:dashboard_result_ch5} thể hiện số liệu tổng hợp sau lần chạy. Hình~\ref{fig:journey_result_ch5} cho thấy một lần mua sắm có đủ bốn khu vực; màu xám là màu quy ước cho nhãn \textit{Neutral}, không phải trạng thái thiếu dữ liệu. Báo cáo chất lượng tại Hình~\ref{fig:quality_result_ch5} tách riêng dữ liệu hợp lệ, ảnh không có khuôn mặt và ảnh không hợp lệ.

\begin{figure}[H]
    \centering
    \includegraphics[width=0.98\textwidth]{Hinhve/Chuong5/5_06_tong_quan_ket_qua.png}
    \caption{Trang tổng quan sau thực nghiệm tích hợp}
    \label{fig:dashboard_result_ch5}
\end{figure}

\begin{figure}[H]
    \centering
    \includegraphics[width=0.98\textwidth]{Hinhve/Chuong5/5_07_hanh_trinh_bon_khu_vuc.png}
    \caption{Lần mua sắm có đủ dữ liệu tại bốn khu vực}
    \label{fig:journey_result_ch5}
\end{figure}

\begin{figure}[H]
    \centering
    \includegraphics[width=0.98\textwidth]{Hinhve/Chuong5/5_08_chat_luong_du_lieu.png}
    \caption{Báo cáo chất lượng dữ liệu sau thực nghiệm}
    \label{fig:quality_result_ch5}
\end{figure}

Các thành phần của hệ thống được đóng gói và khởi động bằng Docker Compose. Giao diện CRM, máy chủ nghiệp vụ, dịch vụ xử lý ảnh và PostgreSQL đều đạt trạng thái sẵn sàng; cổng 8001 của dịch vụ xử lý ảnh chỉ được mở trong môi trường thực nghiệm để kiểm tra trực tiếp. Chín kiểm thử máy chủ nghiệp vụ và ba kiểm thử dịch vụ xử lý ảnh đều đạt; giao diện React và TypeScript biên dịch thành công. Hình~\ref{fig:deployment_status_ch5} và Hình~\ref{fig:test_results_ch5} tổng hợp các bằng chứng này.

\begin{figure}[H]
    \centering
    \includegraphics[width=0.98\textwidth]{Hinhve/Chuong4/4_17_trang_thai_trien_khai.png}
    \caption{Trạng thái của hệ thống tích hợp và cấu hình xử lý ảnh}
    \label{fig:deployment_status_ch5}
\end{figure}

\begin{figure}[H]
    \centering
    \includegraphics[width=0.92\textwidth]{Hinhve/Chuong4/4_18_ket_qua_kiem_thu.png}
    \caption{Kết quả kiểm thử mã nguồn và phát lại dữ liệu ảnh}
    \label{fig:test_results_ch5}
\end{figure}

\subsection{Giới hạn của thực nghiệm}
\label{subsection:5.2.6}

Thứ nhất, thực nghiệm được thực hiện khi chưa có hệ thống camera lắp đặt tại cửa hàng. Các khung quan sát được xây dựng từ ảnh FairFace bằng cách thay đổi vị trí, kích thước và độ sáng, nên không chứa đầy đủ các yếu tố của video thực như chuyển động, thay đổi góc nhìn, mất nét, nhiễu cảm biến, che khuất kéo dài và điều kiện chiếu sáng phức tạp. Khung có hai khuôn mặt chỉ xác nhận rằng một sự kiện có thể tạo nhiều quan sát; trường hợp này chưa đại diện cho mật độ người và mức che khuất trong một không gian mua sắm đông người.

Thứ hai, FairFace không phải tập dữ liệu định danh gồm nhiều ảnh của cùng một người. Trong thực nghiệm, mỗi bản ghi được quy ước thành một khách hàng và các khung đăng ký, quan sát đều được xây dựng từ cùng một ảnh nguồn. Vì vậy, kết quả 115/115 chỉ xác nhận luồng đăng ký, tạo véc-tơ, đối sánh và lưu dữ liệu với đầu vào có kiểm soát. Kết quả không phải là tỷ lệ chính xác nhận dạng khi một người xuất hiện ở ngày khác, góc nhìn khác hoặc trên camera khác.

Thứ ba, ảnh FairFace được sử dụng không có nhãn biểu cảm chuẩn để đối chiếu với kết quả DeepFace. Thực nghiệm chỉ xác nhận rằng mô hình trả được một trong bảy nhãn, mức tin cậy và trạng thái được lưu đúng. Các nhãn được gán theo quy tắc trong bộ dữ liệu kiểm tra nghiệp vụ chỉ phục vụ kiểm tra phép lọc, tổng hợp và hiển thị; chúng không được sử dụng để kết luận về độ chính xác phân loại, trạng thái cảm xúc hoặc mức độ hài lòng của khách hàng.

Thứ tư, 125 lần mua sắm và 90 đơn hàng được xây dựng theo kịch bản vì đề tài chưa kết nối với hệ thống bán hàng thực tế. Do đó, số lượt qua khu vực, sản phẩm được chọn, số lượng mua và mối liên hệ giữa quan sát với đơn hàng chỉ dùng để kiểm tra quy tắc nghiệp vụ. Các phân bố thu được không phản ánh hành vi tiêu dùng của một quần thể khách hàng thực tế.

Thứ năm, trong phạm vi nghiên cứu hiện tại, hệ thống được kiểm nghiệm với dữ liệu của một cửa hàng, chưa theo dõi cùng một người chưa xác định qua nhiều khung hình liên tiếp và chưa xử lý trường hợp một người đồng thời xuất hiện tại nhiều camera. Chỉ năm khách hàng có chuỗi ảnh đầy đủ tại bốn khu vực; 95 khách hàng còn lại được dùng để kiểm tra chức năng đăng ký và nhận dạng tại một khu vực. Các giới hạn này xác định phạm vi của kết luận: thực nghiệm chứng minh các thành phần phần mềm có thể phối hợp trong điều kiện đã xây dựng, chưa chứng minh hiệu quả vận hành trong cửa hàng thực tế. Các giới hạn cấu hình được tổng hợp tại Hình~\ref{fig:version_limits_ch5}.

\begin{figure}[H]
    \centering
    \includegraphics[width=0.92\textwidth]{Hinhve/Chuong4/4_19_gioi_han_phien_ban.png}
    \caption{Phạm vi cấu hình và giới hạn khi diễn giải kết quả}
    \label{fig:version_limits_ch5}
\end{figure}

Báo cáo của hệ thống chỉ mô tả các nhãn biểu cảm được lưu trong phạm vi dữ liệu thử nghiệm; kết quả không được quy đổi thành điểm tích cực--tiêu cực, biểu cảm trung bình hoặc mức độ hài lòng.

\end{document}
